\chapter{Reprodutibilidade}
\label{ap:reprodutibilidade}

Este apêndice documenta a proveniência dos resultados apresentados no corpo da dissertação. Enumera
os artefactos conservados, estabelece a correspondência entre cada resultado reportado e o
procedimento que o produz, apresenta o estado em que ficou cada afirmação substantiva depois de
medida, e regista as condições de licenciamento dos dados utilizados. Encerra com o desenho do
estudo com utilizadores que a Secção~\ref{sec:con_limitacoes} identifica como a lacuna de maior
peso do trabalho.

\section{Artefactos conservados}
\label{sec:ap_artefactos}

Esta secção enumera o que é conservado no repositório de código e o que não é, com a razão de cada
caso.

O ambiente de execução está descrito na Secção~\ref{sec:met_ferramentas}, com as versões exatas das
bibliotecas que produziram os valores deste documento. As versões estão fixadas, e não expressas em
intervalos, porque uma biblioteca que se atualize sozinha altera resultados sem emitir aviso.
Nenhum dos procedimentos exige unidade de processamento gráfico.

São conservadas três classes de artefacto. A primeira é o componente aprendido, que compreende o
modelo treinado, o calibrador ajustado sobre o bloco de validação e um ficheiro descritivo da
execução que os produziu, com a dimensão de cada bloco, a prevalência de casos positivos, a semente
utilizada e as métricas obtidas. A segunda é o conjunto de ficheiros de resultados, um por
procedimento de avaliação, que o texto cita. A terceira é a base de precedentes que a aplicação
implantada consulta, com cinquenta e seis megabytes de vetores e oito de metadados. Esta última não
constitui uma amostra: é o artefacto de que o sistema necessita para arrancar, e sem ele não
existiriam precedentes para apresentar.

Os corpora de origem não são conservados no repositório. O conjunto histórico ocupa várias centenas
de megabytes e é reconstruído pelos procedimentos de preparação descritos na
Secção~\ref{sec:met_dados}. Acompanham o código amostras reduzidas, cuja finalidade é permitir que a
verificação automática decorra sem acesso à rede, incluindo o teste que impõe a separação temporal
descrita na Secção~\ref{sec:met_triagem}.

\section{Origem de cada resultado}
\label{sec:ap_numeros}

Esta secção estabelece a correspondência entre cada resultado reportado no corpo do documento e o
procedimento que o gera.

Nenhum resultado desta dissertação foi introduzido manualmente no texto. Cada um é produzido por um
procedimento automático de semente fixa, que escreve o seu próprio ficheiro de resultados, e é esse
ficheiro que o texto cita. Uma segunda execução sobre os mesmos dados devolve os mesmos valores. A
Tabela~\ref{tab:ap_numeros} reúne os resultados reportados e identifica a secção em que cada um
aparece; o repositório de código nomeia, para cada resultado, o ficheiro que o contém e o
procedimento que o produz.

A garantia de determinismo aplica-se ao bloco superior da tabela. As linhas do bloco inferior são de
outra natureza, e por isso estão separadas: constituem medições sobre a operação do sistema, cujo
valor depende do momento em que o registo foi lido e não de uma semente. A secção em que cada uma é
apresentada declara o período e o conjunto sobre que foi realizada.

Uma dessas medições não é regenerável, e a distinção fica registada em vez de omitida. A repartição
do funil por porta, apresentada na Secção~\ref{sec:sis_politica}, foi lida do registo de decisões
num dia concreto, e esse registo conserva apenas três dias, pela razão indicada na
Secção~\ref{sec:met_dados}: é republicado a cada ciclo, pelo que o custo é de publicação e não de
armazenamento. A execução do mesmo procedimento numa data posterior produz outro dia, e não o mesmo.
O valor consta de um ficheiro de resultados, com a data da leitura e o procedimento que a produziu,
de modo que qualquer leitura futura tenha origem declarada ainda que aquele dia não se repita.

\begin{table}[!htbp]
\caption[Cada resultado e o procedimento que o produz]{Os resultados reportados no corpo do
documento e a secção em que aparecem. O bloco superior é determinístico e reproduz-se exatamente
sobre os mesmos dados; o bloco inferior reúne medições sobre a operação, cujo valor depende da data
de leitura do registo.}
\label{tab:ap_numeros}
\centering
\small
\begin{tabular}{p{0.46\textwidth} L{0.21\textwidth} L{0.19\textwidth}}
\toprule
\tabhead{Resultado} & \tabhead{Valor} & \tabhead{Onde aparece} \\
\midrule
Amplitude da taxa de disparo, desvio relativo contra limiar fixo & $0{,}015$ contra $0{,}344$ &
§\ref{sec:av_qi1} \\
$F_1$ contra o critério aproximado & $0{,}516$ contra $0{,}218$ & §\ref{sec:av_qi1} \\
Comparação com dois detetores aprendidos & $0{,}530$ / $0{,}269$ / $0{,}280$ & §\ref{sec:av_qi1} \\
Ablação à janela e ao peso do desvio & $0{,}678$ e $0{,}664$ contra $0{,}516$ & §\ref{sec:av_qi1} \\
\addlinespace
Precisão@5 da recuperação, quatro métodos & $0{,}514$ / $0{,}346$ / $0{,}240$ / $0{,}126$ &
§\ref{sec:av_qi2} \\
Valor de referência do setor mais representado & $0{,}467$ & §\ref{sec:av_qi2} \\
Precisão@5 no protocolo simétrico à escala & $0{,}595$ com acaso $0{,}333$ & §\ref{sec:av_qi2} \\
Precisão@5 apenas com precedentes anteriores & $0{,}513$ com acaso $0{,}259$ & §\ref{sec:av_qi2} \\
Concordância de direção contra o valor de acaso & $0{,}708$ contra $0{,}688$ & §\ref{sec:av_qi2} \\
\addlinespace
Triagem, \gls{PRAUC} das seis famílias & de $0{,}378$ a $0{,}542$ & §\ref{sec:av_qi3} \\
Precisão dentro do orçamento, quatro políticas & $0{,}163$ / $0{,}379$ / $0{,}662$ / $0{,}632$ &
§\ref{sec:av_qi3} \\
Grelha de nove definições de critério & $9$ de $9$ & §\ref{sec:av_qi3} \\
Ablação da identidade, tabela de consulta contra modelo & $0{,}534$ contra $0{,}538$ &
§\ref{sec:av_ablacao} \\
Acréscimo do texto sobre a tabela de consulta & $0{,}547$ contra $0{,}534$ & §\ref{sec:av_ablacao} \\
\addlinespace
Qualidade do ajuste da decomposição, mediana e casos negativos & $0{,}460$, $1$ em $17$ &
§\ref{sec:av_decomposicao} \\
Seletividade da porta, dentro e entre empresas & $0{,}072$ contra $0{,}392$ &
§\ref{sec:av_producao} \\
Deriva de distribuição, \gls{PSI} da volatilidade & $0{,}281$ & §\ref{sec:av_producao} \\
Sistema completo contra as alternativas & $0{,}489$ / $0{,}632$ / $0{,}662$ / $0{,}968$ &
§\ref{sec:av_producao} \\
\midrule
\multicolumn{3}{l}{\emph{Medições sobre a operação: o valor depende da data de leitura do registo}}
\\
Decisões maturadas, mantidas contra suprimidas & $58{,}9\%$ contra $61{,}7\%$ &
§\ref{sec:av_producao} \\
Capacidade de ordenação sobre as decisões tomadas & $0{,}486$ & §\ref{sec:av_producao} \\
Cobertura de notícias em dias invulgares & $88{,}5\%$ & §\ref{sec:av_producao} \\
Latência, descoberta e entrega & $353$ min e $5$ s & §\ref{sec:sis_operacao} \\
Funil de produção, relevantes até entregues & $944 \rightarrow 42$ & §\ref{sec:sis_politica} \\
Funil de um dia, repartido por porta & $5\,060$ avaliações & §\ref{sec:sis_politica} \\
\bottomrule
\end{tabular}
\end{table}

Os restantes valores numéricos do documento pertencem a três categorias. Os valores extraídos da
literatura estão citados no local em que aparecem. As constantes de configuração constam da
Tabela~\ref{tab:sis_portas}, com a distinção entre as que foram derivadas de uma medição e as que
foram escolhidas. A terceira categoria reúne medições que servem uma passagem e não um resultado,
como as linhas descartadas pelo embargo, o desvio-padrão do caso trabalhado da
Secção~\ref{sec:met_detecao} ou as contagens por fonte de notícias; cada uma declara, no local em
que aparece, o conjunto sobre que foi realizada.

\section{Matriz de evidência}
\label{sec:ap_matriz}

Esta secção submete cada afirmação substantiva do trabalho à evidência que a sustenta e regista o
estado em que ficou.

A tabela da secção anterior enumera resultados. A Tabela~\ref{tab:ap_matriz} enumera afirmações,
exame mais exigente, uma vez que uma afirmação pode não sobreviver à evidência que a suportava. As
afirmações retiradas permanecem na tabela, dado que uma matriz que enumere apenas as sobreviventes
não constitui uma auditoria.

\begin{table}[!htbp]
\caption[Matriz de evidência]{Cada afirmação substantiva, a evidência que a sustenta e o estado em
que ficou. O estado \emph{retirada} indica que o trabalho a abandonou depois de a medir; o estado
\emph{estreitada} indica que subsiste numa formulação mais restrita do que a inicial.}
\label{tab:ap_matriz}
\centering
\scriptsize
\renewcommand{\arraystretch}{1.08}
\begin{tabular}{p{0.34\textwidth} p{0.36\textwidth} L{0.18\textwidth}}
\toprule
\tabhead{Afirmação} & \tabhead{Evidência} & \tabhead{Estado} \\
\midrule
\multicolumn{3}{l}{\textbf{Deteção de movimentos invulgares}} \\
Dispara a ritmo comparável entre empresas de volatilidade muito distinta &
  Amplitude de $0{,}015$ contra $0{,}344$, sem recurso a rótulos & Mantida \\
Supera um limiar fixo no critério aproximado & $F_1$ de $0{,}516$ contra $0{,}218$ &
  Mantida, com o critério declarado como aproximação \\
Supera dois detetores aprendidos com a mesma informação &
  $F_1$ de $0{,}530$ contra $0{,}269$ e $0{,}280$ & Mantida \\
Não utiliza informação posterior ao dia avaliado &
  A janela exclui esse dia, e um teste automático impõe a exclusão & Mantida \\
A janela de vinte dias e o desvio de pesos iguais são as melhores opções &
  Variantes alternativas obtêm $0{,}678$ e $0{,}664$ contra $0{,}516$ &
  \textbf{Retirada}: mantidas por explicabilidade, não por desempenho \\
\midrule
\multicolumn{3}{l}{\textbf{Recuperação de precedentes}} \\
Supera as alternativas lexical e triviais &
  Precisão@5 de $0{,}514$ contra $0{,}346$, $0{,}240$ e $0{,}126$ & Mantida \\
Supera por larga margem o melhor valor de referência trivial &
  A escolha do setor mais representado obtém $0{,}467$, o que reduz a margem de $+0{,}274$ para
  $+0{,}047$ & \textbf{Estreitada}: a afirmação passa a ser por setor, onde se mantém nos cinco \\
Resiste ao aumento de escala do corpus &
  Precisão@5 de $0{,}595$ com acaso de $0{,}333$ & Mantida como teste de escala \\
Resiste à restrição temporal da produção &
  Precisão@5 de $0{,}513$ com acaso de $0{,}259$: a margem varia $0{,}008$ &
  \textbf{Estreitada}: o valor simétrico descreve tarefa menos exigente \\
Capta tema e não direção & Concordância de $0{,}708$ contra acaso de $0{,}688$ &
  Mantida, e reportada como limitação \\
A pertença ao mesmo setor equivale a relação entre notícias & --- &
  \textbf{Não afirmado}: é uma aproximação declarada \\
\midrule
\multicolumn{3}{l}{\textbf{Triagem de notícias}} \\
Nenhuma variante com texto supera a volatilidade &
  \gls{PRAUC} de $0{,}496$ contra $0{,}542$, com nove definições de critério e reamostragem por
  grupos & Mantida \\
O texto acrescenta informação sobre a tabela de consulta &
  \gls{PRAUC} de $0{,}547$ contra $0{,}534$, com intervalo entre $+0{,}004$ e $+0{,}020$ &
  \textbf{Nova}, positiva, e sem efeito na precisão dentro do orçamento \\
As entradas não incorporam informação futura &
  Alterar o futuro de um exemplo deixa as entradas iguais e altera o critério & Mantida \\
A triagem quadruplica a precisão dentro do orçamento &
  O valor de referência de $0{,}163$ resulta de ordenação alfabética; a escolha aleatória obtém
  $0{,}379$ & \textbf{Retirada}: o ganho é de $1{,}67$ vezes \\
O ganho dentro do orçamento provém da aprendizagem &
  Treze constantes por empresa obtêm $0{,}662$ contra $0{,}632$ do modelo & \textbf{Retirada} \\
O modelo implantado seleciona melhor do que o acaso &
  $58{,}9\%$ de casos materiais entre as mantidas contra $61{,}7\%$ entre as suprimidas &
  \textbf{Retirada} \\
\midrule
\multicolumn{3}{l}{\textbf{Sistema}} \\
O texto do alerta não pode divergir do cálculo &
  É composto a partir dos objetos calculados, sem geração de linguagem & Mantida \\
A restrição de gratuitidade limita a cobertura &
  Pelo menos uma notícia relevante em $88{,}5\%$ dos dias invulgares & Mantida como limite superior \\
As explicações são úteis a uma pessoa & --- & \textbf{Não afirmado} \\
\bottomrule
\end{tabular}
\end{table}

Duas observações incidem sobre a tabela no seu conjunto e não sobre qualquer linha isolada.

A primeira é que as afirmações retiradas o foram por medições do próprio trabalho, e nenhuma por
revisão externa. Constitui o comportamento pretendido de um plano de avaliação fixado antes das
experiências, cuja utilidade reside precisamente na capacidade de produzir uma resposta negativa.

A segunda é que as duas linhas assinaladas como não afirmadas têm o mesmo peso das restantes. Uma
delas é a utilidade das explicações para o utilizador, que não foi medida por ausência de estudo com
pessoas. Onde não existe evidência, não é formulada afirmação.

\section{Dados, licenças e atribuição}
\label{sec:ap_dados}

Esta secção regista a origem dos dados utilizados e as obrigações que o seu licenciamento impõe.

O conjunto histórico que sustenta a base de precedentes é o \gls{FNSPID}
\autocite{dong2024fnspid}, utilizado ao abrigo da licença Creative Commons BY-SA 4.0, cuja
atribuição é aqui cumprida. Do corpus constam no documento apenas excertos mínimos de títulos, a
título ilustrativo.

Essa licença comporta uma segunda obrigação, que é acionada neste caso. Uma vez que o repositório
distribui ficheiros derivados do \gls{FNSPID}, a cláusula de partilha nos mesmos termos aplica-se a
esses ficheiros. A licença do código constitui uma decisão que permanece em aberto neste documento,
e é esta a restrição que a condiciona.

Os preços provêm de serviços gratuitos de dados de mercado, através da cadeia de alternativas
descrita na Secção~\ref{sec:sis_dados}; as notícias em tempo real provêm de três interfaces
gratuitas, dentro dos respetivos limites de utilização, comparadas na Secção~\ref{sec:sis_fontes}.

Quanto a dados pessoais, a formulação exata consta da Secção~\ref{sec:met_privacidade}: o sistema
não conhece as posições dos utilizadores e não as solicita, e o canal público não identifica
destinatários. A exceção é a subscrição opcional de alertas personalizados, em que ficam associados
o identificador de conversa atribuído pela plataforma e a lista de empresas escolhida, e esses são
os únicos elementos armazenados.

\section{Desenho do estudo com utilizadores}
\label{sec:ap_estudo}

Esta secção apresenta o desenho do estudo referido na Secção~\ref{sec:con_futuro}.

O estudo compara duas condições sobre seis alertas reais extraídos do canal. A condição A apresenta
o facto sem explicação e a condição B apresenta o alerta completo. Dois dos seis casos são situações
em que o tema da notícia e a direção do movimento discordam, configuração em que a
Secção~\ref{sec:av_qi2} mostra o sistema ser mais suscetível de interpretação incorreta.

O contrabalanço cruza dois fatores, que são a ordem das condições e a metade dos estímulos que serve
de condição A. O segundo fator é necessário porque as duas condições recorrem a alertas distintos,
sob pena de a segunda medir memória em vez de compreensão; sem o cruzamento, qualquer diferença
entre condições ficaria confundida com a maior acessibilidade de uma das metades.

A pergunta central é a que a verificação automática não alcança. A correspondência entre cada frase
do alerta e o facto que a sustenta está garantida por construção, mas a afirmação do produto
descreve um percurso realizado por um leitor: dada uma frase com remissão, aceder ao facto
correspondente e julgar autonomamente se este a sustenta. A questão é qualitativa e não exige
amostra de grande dimensão.

Duas salvaguardas metodológicas estão fixadas no código antes de existirem dados: o limiar de oito
participantes para o teste estatístico, cuja alteração ficaria registada no histórico de versões, e
o procedimento de análise, que sobre a folha por preencher não produz resultado algum. Um terceiro
bloco do desenho original, relativo a texto gerado, não é executável, uma vez que as vias que o
serviam foram retiradas do sistema.
